%! TeX root = ../../main.tex
\documentclass[tikz,border=8pt,11pt]{standalone}
\usepackage[T1]{fontenc}
\usepackage{newtxtext,newtxmath}
\usepackage{xcolor}
\usepackage{tikz}
\usetikzlibrary{shapes.geometric, arrows.meta, positioning, fit, backgrounds, calc}

\begin{document}

% ==============================================================================
% DEFINIZIONE COLORI (PALETTE ACCADEMICA UNIFICATA)
% ==============================================================================
\definecolor{actA}{RGB}{252, 232, 131}          % A: Giallo caldo pastello
\definecolor{actB}{RGB}{242, 168, 178}          % B: Rosa corallo tenue (Azione non-atomica)
\definecolor{actC}{RGB}{188, 222, 178}          % C: Verde salvia
\definecolor{actD}{RGB}{164, 198, 238}          % D: Azzurro polvere (Azione non-atomica)
\definecolor{actNaN}{RGB}{235, 235, 235}        % NaN: Grigio neutro
\definecolor{actEmpty}{RGB}{255, 255, 255}      % Slot vuoto

\definecolor{hdr_bg}{RGB}{242, 245, 249}        % Sfondo intestazioni griglia
\definecolor{hdr_txt}{RGB}{25, 35, 50}          % Testo intestazioni griglia

\begin{tikzpicture}[
    font=\footnotesize,
    x=0.96cm,
    y=0.52cm,
    % Stile cella della griglia
    cell/.style={
        draw=black!80,
        line width=0.55pt,
        inner sep=0pt,
        minimum width=0.96cm,
        minimum height=0.52cm,
        font=\footnotesize,
        text=black,
        align=center
    },
    % Intestazione riga sinistra
    rowhdr/.style={
        draw=black!80,
        line width=0.55pt,
        fill=hdr_bg,
        inner sep=0pt,
        minimum width=2.9cm,
        minimum height=0.52cm,
        font=\footnotesize\bfseries,
        text=hdr_txt,
        align=left
    },
    % Intestazione colonna numerica (1..10)
    colhdr/.style={
        draw=black!80,
        line width=0.55pt,
        fill=hdr_bg,
        inner sep=0pt,
        minimum width=0.96cm,
        minimum height=0.52cm,
        font=\footnotesize\bfseries,
        text=hdr_txt,
        align=center
    },
    % Titolo sopra ogni tabella
    sectiontitle/.style={
        font=\bfseries\footnotesize,
        text=black,
        anchor=south west,
        align=left
    }
]

    % Coordinata X per allineamento a sinistra (bordo sinistro di rowhdr a x = -3.45)
    \def\leftX{-3.45}
    % Coordinata Y per il distacco del titolo dalla tabella
    \def\titleY{0.50}

    % ==========================================================================
    % 1. TRACCE COMPLETE DELL'EVENT LOG
    % ==========================================================================
    \node[sectiontitle] at (\leftX, \titleY) {Tracce complete dell'Event Log};

    % Header colonne
    \node[rowhdr] at (-2.0, 0) {\hspace{6pt}Numero evento};
    \foreach \i in {1,...,10} {
        \node[colhdr] at (\i - 1, 0) {\i};
    }

    % Trace 1
    \node[rowhdr] at (-2.0, -1) {\hspace{6pt}Traccia 1};
    \node[cell, fill=actA] at (0, -1) {A};
    \node[cell, fill=actB] at (1, -1) {B};
    \node[cell, fill=actA] at (2, -1) {A};
    \node[cell, fill=actC] at (3, -1) {C};
    \node[cell, fill=actD] at (4, -1) {D};
    \node[cell, fill=actB] at (5, -1) {B};
    \node[cell, fill=actA] at (6, -1) {A};
    \node[cell, fill=actD] at (7, -1) {D};
    \node[cell, fill=actEmpty] at (8, -1) {};
    \node[cell, fill=actEmpty] at (9, -1) {};

    % Trace 2
    \node[rowhdr] at (-2.0, -2) {\hspace{6pt}Traccia 2};
    \node[cell, fill=actA] at (0, -2) {A};
    \node[cell, fill=actC] at (1, -2) {C};
    \node[cell, fill=actA] at (2, -2) {A};
    \node[cell, fill=actA] at (3, -2) {A};
    \node[cell, fill=actB] at (4, -2) {B};
    \node[cell, fill=actD] at (5, -2) {D};
    \node[cell, fill=actD] at (6, -2) {D};
    \node[cell, fill=actC] at (7, -2) {C};
    \node[cell, fill=actA] at (8, -2) {A};
    \node[cell, fill=actA] at (9, -2) {A};

    % ==========================================================================
    % 2. PREFIX_LENGTH, lunghezza = 5
    % ==========================================================================
    \begin{scope}[yshift=-2.85cm]
        \node[sectiontitle] at (\leftX, \titleY) {\texttt{PREFIX\_LENGTH}, lunghezza $= 5$};

        % Header colonne
        \node[rowhdr] at (-2.0, 0) {\hspace{6pt}Numero evento};
        \foreach \i in {1,...,10} {
            \node[colhdr] at (\i - 1, 0) {\i};
        }

        % Prefix 1
        \node[rowhdr] at (-2.0, -1) {\hspace{6pt}Prefisso 1};
        \node[cell, fill=actA] at (0, -1) {A};
        \node[cell, fill=actB] at (1, -1) {B};
        \node[cell, fill=actA] at (2, -1) {A};
        \node[cell, fill=actC] at (3, -1) {C};
        \node[cell, fill=actD] at (4, -1) {D};
        \foreach \i in {5,...,9} {
            \node[cell, fill=actEmpty] at (\i, -1) {};
        }

        % Prefix 2
        \node[rowhdr] at (-2.0, -2) {\hspace{6pt}Prefisso 2};
        \node[cell, fill=actA] at (0, -2) {A};
        \node[cell, fill=actC] at (1, -2) {C};
        \node[cell, fill=actA] at (2, -2) {A};
        \node[cell, fill=actA] at (3, -2) {A};
        \node[cell, fill=actB] at (4, -2) {B};
        \foreach \i in {5,...,9} {
            \node[cell, fill=actEmpty] at (\i, -2) {};
        }
    \end{scope}

    % ==========================================================================
    % 3. COMPONENT, target = [A,B,C], componente = 2
    % ==========================================================================
    \begin{scope}[yshift=-5.70cm]
        \node[sectiontitle] at (\leftX, \titleY) {\texttt{COMPONENT}, target $= [\text{A},\text{B},\text{C}]$, componente $= 2$};

        % Header colonne
        \node[rowhdr] at (-2.0, 0) {\hspace{6pt}Numero evento};
        \foreach \i in {1,...,10} {
            \node[colhdr] at (\i - 1, 0) {\i};
        }

        % Prefix 1
        \node[rowhdr] at (-2.0, -1) {\hspace{6pt}Prefisso 1};
        \node[cell, fill=actA] at (0, -1) {A};
        \node[cell, fill=actB] at (1, -1) {B};
        \node[cell, fill=actA] at (2, -1) {A};
        \node[cell, fill=actC] at (3, -1) {C};
        \node[cell, fill=actD] at (4, -1) {D};
        \node[cell, fill=actB] at (5, -1) {B};
        \node[cell, fill=actNaN] at (6, -1) {\footnotesize\textit{NaN}};
        \node[cell, fill=actNaN] at (7, -1) {\footnotesize\textit{NaN}};
        \node[cell, fill=actEmpty] at (8, -1) {};
        \node[cell, fill=actEmpty] at (9, -1) {};

        % Prefix 2
        \node[rowhdr] at (-2.0, -2) {\hspace{6pt}Prefisso 2};
        \node[cell, fill=actA] at (0, -2) {A};
        \node[cell, fill=actC] at (1, -2) {C};
        \node[cell, fill=actA] at (2, -2) {A};
        \node[cell, fill=actA] at (3, -2) {A};
        \node[cell, fill=actB] at (4, -2) {B};
        \node[cell, fill=actD] at (5, -2) {D};
        \node[cell, fill=actD] at (6, -2) {D};
        \node[cell, fill=actC] at (7, -2) {C};
        \node[cell, fill=actEmpty] at (8, -2) {};
        \node[cell, fill=actEmpty] at (9, -2) {};
    \end{scope}

    % ==========================================================================
    % 4. COMPONENT_TIME_WARPED
    % ==========================================================================
    \begin{scope}[yshift=-8.55cm]
        \node[sectiontitle] at (\leftX, \titleY) {\texttt{COMPONENT\_TIME\_WARPED}, target $= [\text{A},\text{B},\text{C}]$, componente $= 2$, non-atomiche $= [\text{B},\text{D}]$};

        % Header colonne
        \node[rowhdr] at (-2.0, 0) {\hspace{6pt}Numero evento};
        \foreach \i in {1,...,10} {
            \node[colhdr] at (\i - 1, 0) {\i};
        }

        % Prefix 1 (D duplicato a pos 5 e 6, B duplicato a pos 7 e 8)
        \node[rowhdr] at (-2.0, -1) {\hspace{6pt}Prefisso 1};
        \node[cell, fill=actA] at (0, -1) {A};
        \node[cell, fill=actB] at (1, -1) {B};
        \node[cell, fill=actA] at (2, -1) {A};
        \node[cell, fill=actC] at (3, -1) {C};
        \node[cell, fill=actD] at (4, -1) {D};
        \node[cell, fill=actD] at (5, -1) {D};
        \node[cell, fill=actB] at (6, -1) {B};
        \node[cell, fill=actB] at (7, -1) {B};
        \node[cell, fill=actEmpty] at (8, -1) {};
        \node[cell, fill=actEmpty] at (9, -1) {};

        % Prefix 2
        \node[rowhdr] at (-2.0, -2) {\hspace{6pt}Prefisso 2};
        \node[cell, fill=actA] at (0, -2) {A};
        \node[cell, fill=actC] at (1, -2) {C};
        \node[cell, fill=actA] at (2, -2) {A};
        \node[cell, fill=actA] at (3, -2) {A};
        \node[cell, fill=actB] at (4, -2) {B};
        \node[cell, fill=actD] at (5, -2) {D};
        \node[cell, fill=actD] at (6, -2) {D};
        \node[cell, fill=actC] at (7, -2) {C};
        \node[cell, fill=actEmpty] at (8, -2) {};
        \node[cell, fill=actEmpty] at (9, -2) {};
    \end{scope}

\end{tikzpicture}

\end{document}
